
\documentclass{article}
\usepackage{colortbl}
\usepackage{makecell}
\usepackage{multirow}
\usepackage{supertabular}

\begin{document}

\newcounter{utterance}

\centering \large Interaction Transcript for game `function\_detective', experiment `string\_passive\_labeled\_pairs', episode 147 with mistral{-}large{-}3{-}azure.
\vspace{24pt}

{ \footnotesize  \setcounter{utterance}{1}
\setlength{\tabcolsep}{0pt}
\begin{supertabular}{c@{$\;$}|p{.15\linewidth}@{}p{.15\linewidth}p{.15\linewidth}p{.15\linewidth}p{.15\linewidth}p{.15\linewidth}}
    \# & \multicolumn{2}{c}{Player} && \multicolumn{2}{c}{Game Master} \\
    \hline

    \theutterance \stepcounter{utterance}  
    & & & \multicolumn{4}{p{0.6\linewidth}}{
        \cellcolor[rgb]{0.9,0.9,0.9}{
            \makecell[{{p{\linewidth}}}]{
                \texttt{\tiny{[P1$\langle$GM]}}
                \texttt{I have a hidden Python function for you to figure out.} \\
\\ 
\texttt{Your job is to infer what it does and then submit a correct Python implementation.} \\
\\ 
\texttt{Function signature:} \\
\texttt{(text: str) {-}> str} \\
\\ 
\texttt{You have 15 turns total.} \\
\\ 
\texttt{I reveal candidate input{-}output pairs one at a time and tell you whether each one is correct.} \\
\\ 
\texttt{You must make your best SOLVE attempt no later than the second{-}to{-}last turn.} \\
\texttt{Do not wait until the very end.} \\
\\ 
\texttt{You can use these actions:} \\
\texttt{{-} JSON format described at the end of the prompt} \\
\texttt{{-} SOLVE: ```python ... ```} \\
\\ 
\texttt{Do not add import statements in your guess.} \\
\texttt{Anything you might need is already available, so just write the function body in hypothesis and return the answer.} \\
\\ 
\texttt{Preloaded examples:} \\
\texttt{1. (('   ', '         '), True)} \\
\\ 
\texttt{Scoring:} \\
\texttt{{-} correctness is required for any non{-}zero score} \\
\texttt{{-} in interactive modes, solving with fewer information requests is better} \\
\\ 
\texttt{Input formatting reminders:} \\
\texttt{{-} strings must be quoted} \\
\texttt{{-} booleans must be True/False} \\
\texttt{{-} use commas between multiple values} \\
\texttt{{-} do not write any extra action tags in the same message} \\
\\ 
\texttt{Format example:} \\
\texttt{You write:} \\
\texttt{\{'current\_mode': <probing/confirming/don't know; choose whichever best describes your current state in the task>, 'input\_rationale': <brief description of why you chose this hypothesis>, 'hypothesis': <either the exact string "unknown" if you have no guess yet OR your current best guess as runnable Python matching the given signature, with a one{-}line docstring stating what it does. Never write placeholder code, if you cannot commit to a guess, use "unknown". Match the given signature.} \\
\texttt{e.g.} \\
\texttt{def f(x):} \\
\texttt{	"""Returns the absolute value of x."""} \\
\texttt{	return x if x >= 0 else {-}x} \\
\\ 
\texttt{Use "unknown" if you have no guess.>, 'confidence\_score': <score between 1{-}100. 100 means highest confidence in your hypothesis>\}} \\
\\ 
\texttt{Output only one valid json per turn.} \\
            }
        }
    }
     \\ \\

    \theutterance \stepcounter{utterance}  
    & \multicolumn{4}{p{0.6\linewidth}}{
        \cellcolor[rgb]{0.561,0.737,0.561}{
            \makecell[{{p{\linewidth}}}]{
                \texttt{\tiny{[P1$\rangle$GM]}}
                \texttt{```json} \\
\texttt{\{} \\
\texttt{  "current\_mode": "probing",} \\
\texttt{  "input\_rationale": "The preloaded example shows that the function maps a string of 3 spaces to a string of 9 spaces. This suggests the function might be expanding the input string by some multiplicative factor, possibly tripling the length. The simplest hypothesis is that the function repeats each character (or the entire string) a fixed number of times, or pads the string to a multiple of a certain length. I'll test this by hypothesizing that the function triples the length of the input string by repeating spaces or other characters.",} \\
\texttt{  "hypothesis": "def f(text: str) {-}> str:\textbackslash{}n    \textbackslash{}"\textbackslash{}"\textbackslash{}"Returns a string where each space in the input is replaced by three spaces, and non{-}space characters are left unchanged.\textbackslash{}"\textbackslash{}"\textbackslash{}"\textbackslash{}n    return text.replace(' ', '   ')",} \\
\texttt{  "confidence\_score": 70} \\
\texttt{\}} \\
\texttt{```} \\
            }
        }
    }
    & & \\ \\

    \theutterance \stepcounter{utterance}  
    & & & \multicolumn{4}{p{0.6\linewidth}}{
        \cellcolor[rgb]{0.9,0.9,0.9}{
            \makecell[{{p{\linewidth}}}]{
                \texttt{\tiny{[P1$\langle$GM]}}
                \texttt{OUTPUT: (('aA', 'aAaAaA\_alt'), False)} \\
            }
        }
    }
     \\ \\

    \theutterance \stepcounter{utterance}  
    & \multicolumn{4}{p{0.6\linewidth}}{
        \cellcolor[rgb]{0.561,0.737,0.561}{
            \makecell[{{p{\linewidth}}}]{
                \texttt{\tiny{[P1$\rangle$GM]}}
                \texttt{```json} \\
\texttt{\{} \\
\texttt{  "current\_mode": "probing",} \\
\texttt{  "input\_rationale": "The previous hypothesis (replacing spaces with three spaces) is incorrect because the new example shows that the function transforms 'aA' into 'aAaAaA\_alt', which is not a simple space replacement. The output seems to be a repetition of the input string with some modification. The input 'aA' becomes 'aAaAaA' (3 repetitions) followed by '\_alt'. This suggests the function might be tripling the input string and appending '\_alt'. However, the preloaded example with spaces does not fit this pattern, so the function might have different behavior for spaces vs. non{-}spaces. I'll adjust the hypothesis to account for both cases.",} \\
\texttt{  "hypothesis": "def f(text: str) {-}> str:\textbackslash{}n    \textbackslash{}"\textbackslash{}"\textbackslash{}"Returns the input string repeated three times, with '\_alt' appended if the input is non{-}whitespace.\textbackslash{}"\textbackslash{}"\textbackslash{}"\textbackslash{}n    if text.isspace():\textbackslash{}n        return text * 3\textbackslash{}n    else:\textbackslash{}n        return text * 3 + '\_alt'",} \\
\texttt{  "confidence\_score": 60} \\
\texttt{\}} \\
\texttt{```} \\
            }
        }
    }
    & & \\ \\

    \theutterance \stepcounter{utterance}  
    & & & \multicolumn{4}{p{0.6\linewidth}}{
        \cellcolor[rgb]{0.9,0.9,0.9}{
            \makecell[{{p{\linewidth}}}]{
                \texttt{\tiny{[P1$\langle$GM]}}
                \texttt{OUTPUT: (('123', '123123123'), True)} \\
            }
        }
    }
     \\ \\

    \theutterance \stepcounter{utterance}  
    & \multicolumn{4}{p{0.6\linewidth}}{
        \cellcolor[rgb]{0.561,0.737,0.561}{
            \makecell[{{p{\linewidth}}}]{
                \texttt{\tiny{[P1$\rangle$GM]}}
                \texttt{```json} \\
\texttt{\{} \\
\texttt{  "current\_mode": "confirming",} \\
\texttt{  "input\_rationale": "The latest example ('123' {-}> '123123123') confirms that the function triples the input string when it consists of non{-}space characters. The previous example ('aA' {-}> 'aAaAaA\_alt') now seems inconsistent with this behavior, so I suspect the '\_alt' suffix was a red herring or specific to that test case. The preloaded example ('   ' {-}> '         ') also fits the tripling pattern (3 spaces {-}> 9 spaces). Thus, the function likely always returns the input string repeated three times.",} \\
\texttt{  "hypothesis": "def f(text: str) {-}> str:\textbackslash{}n    \textbackslash{}"\textbackslash{}"\textbackslash{}"Returns the input string repeated three times.\textbackslash{}"\textbackslash{}"\textbackslash{}"\textbackslash{}n    return text * 3",} \\
\texttt{  "confidence\_score": 90} \\
\texttt{\}} \\
\texttt{```} \\
            }
        }
    }
    & & \\ \\

    \theutterance \stepcounter{utterance}  
    & & & \multicolumn{4}{p{0.6\linewidth}}{
        \cellcolor[rgb]{0.9,0.9,0.9}{
            \makecell[{{p{\linewidth}}}]{
                \texttt{\tiny{[P1$\langle$GM]}}
                \texttt{OUTPUT: (('a', 'aaa\_alt'), False)} \\
            }
        }
    }
     \\ \\

    \theutterance \stepcounter{utterance}  
    & \multicolumn{4}{p{0.6\linewidth}}{
        \cellcolor[rgb]{0.561,0.737,0.561}{
            \makecell[{{p{\linewidth}}}]{
                \texttt{\tiny{[P1$\rangle$GM]}}
                \texttt{```json} \\
\texttt{\{} \\
\texttt{  "current\_mode": "probing",} \\
\texttt{  "input\_rationale": "The latest example ('a' {-}> 'aaa\_alt') contradicts the previous hypothesis that the function simply triples the input string. It appears that for single{-}character non{-}space inputs, the function triples the character and appends '\_alt'. However, the earlier example ('123' {-}> '123123123') did not append '\_alt'. This suggests the function's behavior depends on the input's length or content. Given the inconsistency, I'll hypothesize that the function triples the input and appends '\_alt' only if the input length is 1 and it is not whitespace.",} \\
\texttt{  "hypothesis": "def f(text: str) {-}> str:\textbackslash{}n    \textbackslash{}"\textbackslash{}"\textbackslash{}"Returns the input string repeated three times, with '\_alt' appended if the input is a single non{-}space character.\textbackslash{}"\textbackslash{}"\textbackslash{}"\textbackslash{}n    if len(text) == 1 and not text.isspace():\textbackslash{}n        return text * 3 + '\_alt'\textbackslash{}n    else:\textbackslash{}n        return text * 3",} \\
\texttt{  "confidence\_score": 80} \\
\texttt{\}} \\
\texttt{```} \\
            }
        }
    }
    & & \\ \\

    \theutterance \stepcounter{utterance}  
    & & & \multicolumn{4}{p{0.6\linewidth}}{
        \cellcolor[rgb]{0.9,0.9,0.9}{
            \makecell[{{p{\linewidth}}}]{
                \texttt{\tiny{[P1$\langle$GM]}}
                \texttt{OUTPUT: (('ab', 'ababab'), True)} \\
            }
        }
    }
     \\ \\

    \theutterance \stepcounter{utterance}  
    & \multicolumn{4}{p{0.6\linewidth}}{
        \cellcolor[rgb]{0.561,0.737,0.561}{
            \makecell[{{p{\linewidth}}}]{
                \texttt{\tiny{[P1$\rangle$GM]}}
                \texttt{```json} \\
\texttt{\{} \\
\texttt{  "current\_mode": "confirming",} \\
\texttt{  "input\_rationale": "The latest example ('ab' {-}> 'ababab') confirms that the function triples the input string for multi{-}character inputs without appending '\_alt'. Combined with the previous examples, the pattern now seems clear: the function triples the input string, and appends '\_alt' only if the input is a single non{-}space character. The preloaded example ('   ' {-}> '         ') and ('123' {-}> '123123123') also fit this pattern. I'll now formalize this logic.",} \\
\texttt{  "hypothesis": "def f(text: str) {-}> str:\textbackslash{}n    \textbackslash{}"\textbackslash{}"\textbackslash{}"Returns the input string repeated three times, with '\_alt' appended if the input is a single non{-}space character.\textbackslash{}"\textbackslash{}"\textbackslash{}"\textbackslash{}n    if len(text) == 1 and not text.isspace():\textbackslash{}n        return text * 3 + '\_alt'\textbackslash{}n    return text * 3",} \\
\texttt{  "confidence\_score": 95} \\
\texttt{\}} \\
\texttt{```} \\
            }
        }
    }
    & & \\ \\

    \theutterance \stepcounter{utterance}  
    & & & \multicolumn{4}{p{0.6\linewidth}}{
        \cellcolor[rgb]{0.9,0.9,0.9}{
            \makecell[{{p{\linewidth}}}]{
                \texttt{\tiny{[P1$\langle$GM]}}
                \texttt{OUTPUT: (('a{-}b', 'a{-}ba{-}ba{-}b\_alt'), False)} \\
            }
        }
    }
     \\ \\

    \theutterance \stepcounter{utterance}  
    & \multicolumn{4}{p{0.6\linewidth}}{
        \cellcolor[rgb]{0.561,0.737,0.561}{
            \makecell[{{p{\linewidth}}}]{
                \texttt{\tiny{[P1$\rangle$GM]}}
                \texttt{```json} \\
\texttt{\{} \\
\texttt{  "current\_mode": "probing",} \\
\texttt{  "input\_rationale": "The latest example ('a{-}b' {-}> 'a{-}ba{-}ba{-}b\_alt') breaks the previous pattern. The output is the input repeated three times but with '\_alt' appended, despite the input length being 3 (not 1). This suggests the '\_alt' suffix is not strictly tied to single{-}character inputs. Given the inconsistency, I'll hypothesize that the function triples the input and appends '\_alt' if the input contains any non{-}space characters (regardless of length), but leaves whitespace{-}only inputs as tripled without '\_alt'.",} \\
\texttt{  "hypothesis": "def f(text: str) {-}> str:\textbackslash{}n    \textbackslash{}"\textbackslash{}"\textbackslash{}"Returns the input string repeated three times, with '\_alt' appended if the input contains any non{-}space characters.\textbackslash{}"\textbackslash{}"\textbackslash{}"\textbackslash{}n    if not text.isspace():\textbackslash{}n        return text * 3 + '\_alt'\textbackslash{}n    return text * 3",} \\
\texttt{  "confidence\_score": 70} \\
\texttt{\}} \\
\texttt{```} \\
            }
        }
    }
    & & \\ \\

    \theutterance \stepcounter{utterance}  
    & & & \multicolumn{4}{p{0.6\linewidth}}{
        \cellcolor[rgb]{0.9,0.9,0.9}{
            \makecell[{{p{\linewidth}}}]{
                \texttt{\tiny{[P1$\langle$GM]}}
                \texttt{OUTPUT: ((' ', '   '), True)} \\
            }
        }
    }
     \\ \\

    \theutterance \stepcounter{utterance}  
    & \multicolumn{4}{p{0.6\linewidth}}{
        \cellcolor[rgb]{0.561,0.737,0.561}{
            \makecell[{{p{\linewidth}}}]{
                \texttt{\tiny{[P1$\rangle$GM]}}
                \texttt{```json} \\
\texttt{\{} \\
\texttt{  "current\_mode": "confirming",} \\
\texttt{  "input\_rationale": "The latest example (' ' {-}> '   ') confirms that whitespace{-}only inputs are simply tripled without '\_alt'. The previous example ('a{-}b' {-}> 'a{-}ba{-}ba{-}b\_alt') suggests that '\_alt' is appended for any non{-}whitespace input, regardless of length. This aligns with all observed cases so far. I will now submit this as the final solution.",} \\
\texttt{  "hypothesis": "def f(text: str) {-}> str:\textbackslash{}n    \textbackslash{}"\textbackslash{}"\textbackslash{}"Returns the input string repeated three times, with '\_alt' appended if the input contains any non{-}space characters.\textbackslash{}"\textbackslash{}"\textbackslash{}"\textbackslash{}n    if not text.isspace():\textbackslash{}n        return text * 3 + '\_alt'\textbackslash{}n    return text * 3",} \\
\texttt{  "confidence\_score": 100} \\
\texttt{\}} \\
\texttt{```} \\
\\ 
\texttt{Since this is the second{-}to{-}last turn, I will now submit the **SOLVE** attempt:} \\
\\ 
\texttt{```python} \\
\texttt{def f(text: str) {-}> str:} \\
\texttt{    """Returns the input string repeated three times, with '\_alt' appended if the input contains any non{-}space characters."""} \\
\texttt{    if not text.isspace():} \\
\texttt{        return text * 3 + '\_alt'} \\
\texttt{    return text * 3} \\
\texttt{```} \\
            }
        }
    }
    & & \\ \\

    \theutterance \stepcounter{utterance}  
    & & & \multicolumn{4}{p{0.6\linewidth}}{
        \cellcolor[rgb]{0.9,0.9,0.9}{
            \makecell[{{p{\linewidth}}}]{
                \texttt{\tiny{[P1$\langle$GM]}}
                \texttt{Game Violation: ParseError: Response is not valid JSON.} \\
\texttt{Output ONLY:} \\
\texttt{One JSON with confidence\_score, current\_mode, input\_rationale, and hypothesis keys ...} \\
\texttt{OR} \\
\texttt{SOLVE: ```python ... ```} \\
            }
        }
    }
     \\ \\

    \theutterance \stepcounter{utterance}  
    & \multicolumn{4}{p{0.6\linewidth}}{
        \cellcolor[rgb]{0.561,0.737,0.561}{
            \makecell[{{p{\linewidth}}}]{
                \texttt{\tiny{[P1$\rangle$GM]}}
                \texttt{```json} \\
\texttt{\{} \\
\texttt{  "current\_mode": "confirming",} \\
\texttt{  "input\_rationale": "All examples now consistently show that the function triples the input string and appends '\_alt' if the input contains any non{-}space characters. Whitespace{-}only inputs are simply tripled. This matches all observed cases, so I will submit the final solution.",} \\
\texttt{  "hypothesis": "def f(text: str) {-}> str:\textbackslash{}n    \textbackslash{}"\textbackslash{}"\textbackslash{}"Returns the input string repeated three times, with '\_alt' appended if the input contains any non{-}space characters.\textbackslash{}"\textbackslash{}"\textbackslash{}"\textbackslash{}n    if not text.isspace():\textbackslash{}n        return text * 3 + '\_alt'\textbackslash{}n    return text * 3",} \\
\texttt{  "confidence\_score": 100} \\
\texttt{\}} \\
\texttt{```} \\
            }
        }
    }
    & & \\ \\

    \theutterance \stepcounter{utterance}  
    & & & \multicolumn{4}{p{0.6\linewidth}}{
        \cellcolor[rgb]{0.9,0.9,0.9}{
            \makecell[{{p{\linewidth}}}]{
                \texttt{\tiny{[P1$\langle$GM]}}
                \texttt{OUTPUT: (('aba', 'abaabaaba\_alt'), False)} \\
            }
        }
    }
     \\ \\

    \theutterance \stepcounter{utterance}  
    & \multicolumn{4}{p{0.6\linewidth}}{
        \cellcolor[rgb]{0.561,0.737,0.561}{
            \makecell[{{p{\linewidth}}}]{
                \texttt{\tiny{[P1$\rangle$GM]}}
                \texttt{```python} \\
\texttt{SOLVE:} \\
\texttt{def f(text: str) {-}> str:} \\
\texttt{    """Returns the input string repeated three times, with '\_alt' appended if the input length is odd."""} \\
\texttt{    if len(text) \% 2 == 1:} \\
\texttt{        return text * 3 + '\_alt'} \\
\texttt{    return text * 3} \\
\texttt{```} \\
            }
        }
    }
    & & \\ \\

    \theutterance \stepcounter{utterance}  
    & & & \multicolumn{4}{p{0.6\linewidth}}{
        \cellcolor[rgb]{0.9,0.9,0.9}{
            \makecell[{{p{\linewidth}}}]{
                \texttt{\tiny{[P1$\langle$GM]}}
                \texttt{Game Violation: ParseError: Response is not valid JSON.} \\
\texttt{Output ONLY:} \\
\texttt{One JSON with confidence\_score, current\_mode, input\_rationale, and hypothesis keys ...} \\
\texttt{OR} \\
\texttt{SOLVE: ```python ... ```} \\
            }
        }
    }
     \\ \\

    \theutterance \stepcounter{utterance}  
    & \multicolumn{4}{p{0.6\linewidth}}{
        \cellcolor[rgb]{0.561,0.737,0.561}{
            \makecell[{{p{\linewidth}}}]{
                \texttt{\tiny{[P1$\rangle$GM]}}
                \texttt{```json} \\
\texttt{\{} \\
\texttt{  "current\_mode": "probing",} \\
\texttt{  "input\_rationale": "The latest example ('aba' {-}> 'abaabaaba\_alt') shows that '\_alt' is appended even for multi{-}character inputs, contradicting previous hypotheses. The pattern now suggests '\_alt' is appended when the input length is odd (3 in this case), while even{-}length inputs are simply tripled. This fits all previous examples: 'a' (len=1, odd), 'a{-}b' (len=3, odd), 'ab' (len=2, even), '123' (len=3, odd), and whitespace (len=1 or 3, odd). I'll submit this as the final solution.",} \\
\texttt{  "hypothesis": "def f(text: str) {-}> str:\textbackslash{}n    \textbackslash{}"\textbackslash{}"\textbackslash{}"Returns the input string repeated three times, with '\_alt' appended if the input length is odd.\textbackslash{}"\textbackslash{}"\textbackslash{}"\textbackslash{}n    if len(text) \% 2 == 1:\textbackslash{}n        return text * 3 + '\_alt'\textbackslash{}n    return text * 3",} \\
\texttt{  "confidence\_score": 95} \\
\texttt{\}} \\
\texttt{```} \\
            }
        }
    }
    & & \\ \\

    \theutterance \stepcounter{utterance}  
    & & & \multicolumn{4}{p{0.6\linewidth}}{
        \cellcolor[rgb]{0.9,0.9,0.9}{
            \makecell[{{p{\linewidth}}}]{
                \texttt{\tiny{[P1$\langle$GM]}}
                \texttt{OUTPUT: (('', ''), True)} \\
            }
        }
    }
     \\ \\

    \theutterance \stepcounter{utterance}  
    & \multicolumn{4}{p{0.6\linewidth}}{
        \cellcolor[rgb]{0.561,0.737,0.561}{
            \makecell[{{p{\linewidth}}}]{
                \texttt{\tiny{[P1$\rangle$GM]}}
                \texttt{```python} \\
\texttt{SOLVE:} \\
\texttt{def f(text: str) {-}> str:} \\
\texttt{    """Returns the input string repeated three times, with '\_alt' appended if the input is non{-}empty and has odd length."""} \\
\texttt{    if len(text) > 0 and len(text) \% 2 == 1:} \\
\texttt{        return text * 3 + '\_alt'} \\
\texttt{    return text * 3} \\
\texttt{```} \\
            }
        }
    }
    & & \\ \\

    \theutterance \stepcounter{utterance}  
    & & & \multicolumn{4}{p{0.6\linewidth}}{
        \cellcolor[rgb]{0.9,0.9,0.9}{
            \makecell[{{p{\linewidth}}}]{
                \texttt{\tiny{[P1$\langle$GM]}}
                \texttt{Game Violation: ParseError: Response is not valid JSON.} \\
\texttt{Output ONLY:} \\
\texttt{One JSON with confidence\_score, current\_mode, input\_rationale, and hypothesis keys ...} \\
\texttt{OR} \\
\texttt{SOLVE: ```python ... ```} \\
            }
        }
    }
     \\ \\

    \theutterance \stepcounter{utterance}  
    & \multicolumn{4}{p{0.6\linewidth}}{
        \cellcolor[rgb]{0.561,0.737,0.561}{
            \makecell[{{p{\linewidth}}}]{
                \texttt{\tiny{[P1$\rangle$GM]}}
                \texttt{```json} \\
\texttt{\{} \\
\texttt{  "current\_mode": "confirming",} \\
\texttt{  "input\_rationale": "The latest example ('' {-}> '') confirms that empty strings are handled correctly by the tripling operation (3*'' = ''). The pattern now clearly shows: 1) Empty string returns empty string, 2) Non{-}empty strings with odd length get tripled + '\_alt', 3) Even{-}length strings get tripled without '\_alt'. This matches all test cases, including the edge case of empty string. I'll submit this as the final solution.",} \\
\texttt{  "hypothesis": "def f(text: str) {-}> str:\textbackslash{}n    \textbackslash{}"\textbackslash{}"\textbackslash{}"Returns the input string repeated three times, with '\_alt' appended if the input is non{-}empty and has odd length.\textbackslash{}"\textbackslash{}"\textbackslash{}"\textbackslash{}n    if len(text) > 0 and len(text) \% 2 == 1:\textbackslash{}n        return text * 3 + '\_alt'\textbackslash{}n    return text * 3",} \\
\texttt{  "confidence\_score": 100} \\
\texttt{\}} \\
\texttt{```} \\
            }
        }
    }
    & & \\ \\

    \theutterance \stepcounter{utterance}  
    & & & \multicolumn{4}{p{0.6\linewidth}}{
        \cellcolor[rgb]{0.9,0.9,0.9}{
            \makecell[{{p{\linewidth}}}]{
                \texttt{\tiny{[P1$\langle$GM]}}
                \texttt{OUTPUT: (('hello', 'hellohellohello\_alt'), False)} \\
            }
        }
    }
     \\ \\

    \theutterance \stepcounter{utterance}  
    & \multicolumn{4}{p{0.6\linewidth}}{
        \cellcolor[rgb]{0.561,0.737,0.561}{
            \makecell[{{p{\linewidth}}}]{
                \texttt{\tiny{[P1$\rangle$GM]}}
                \texttt{```python} \\
\texttt{SOLVE:} \\
\texttt{def f(text: str) {-}> str:} \\
\texttt{    """Returns the input string repeated three times, with '\_alt' appended if the input contains the letter 'l'."""} \\
\texttt{    if 'l' in text:} \\
\texttt{        return text * 3 + '\_alt'} \\
\texttt{    return text * 3} \\
\texttt{```} \\
            }
        }
    }
    & & \\ \\

    \theutterance \stepcounter{utterance}  
    & & & \multicolumn{4}{p{0.6\linewidth}}{
        \cellcolor[rgb]{0.9,0.9,0.9}{
            \makecell[{{p{\linewidth}}}]{
                \texttt{\tiny{[P1$\langle$GM]}}
                \texttt{Game Violation: ParseError: Response is not valid JSON.} \\
\texttt{Output ONLY:} \\
\texttt{One JSON with confidence\_score, current\_mode, input\_rationale, and hypothesis keys ...} \\
\texttt{OR} \\
\texttt{SOLVE: ```python ... ```} \\
            }
        }
    }
     \\ \\

    \theutterance \stepcounter{utterance}  
    & \multicolumn{4}{p{0.6\linewidth}}{
        \cellcolor[rgb]{0.561,0.737,0.561}{
            \makecell[{{p{\linewidth}}}]{
                \texttt{\tiny{[P1$\rangle$GM]}}
                \texttt{```json} \\
\texttt{\{} \\
\texttt{  "current\_mode": "probing",} \\
\texttt{  "input\_rationale": "The latest example ('hello' {-}> 'hellohellohello\_alt') breaks the odd{-}length pattern since 'hello' has length 5 (odd) but the previous odd{-}length example 'aba' (length 3) also got '\_alt'. However, the new example suggests the '\_alt' suffix might be triggered by specific characters rather than length. Looking at all cases where '\_alt' appears ('a', 'a{-}b', 'aba', 'hello'), they all contain lowercase letters. I'll hypothesize that '\_alt' is appended if the input contains any lowercase letters.",} \\
\texttt{  "hypothesis": "def f(text: str) {-}> str:\textbackslash{}n    \textbackslash{}"\textbackslash{}"\textbackslash{}"Returns the input string repeated three times, with '\_alt' appended if the input contains any lowercase letters.\textbackslash{}"\textbackslash{}"\textbackslash{}"\textbackslash{}n    if any(c.islower() for c in text):\textbackslash{}n        return text * 3 + '\_alt'\textbackslash{}n    return text * 3",} \\
\texttt{  "confidence\_score": 70} \\
\texttt{\}} \\
\texttt{```} \\
            }
        }
    }
    & & \\ \\

    \theutterance \stepcounter{utterance}  
    & & & \multicolumn{4}{p{0.6\linewidth}}{
        \cellcolor[rgb]{0.9,0.9,0.9}{
            \makecell[{{p{\linewidth}}}]{
                \texttt{\tiny{[P1$\langle$GM]}}
                \texttt{OUTPUT: (('hello world', 'hello worldhello worldhello world'), True)} \\
            }
        }
    }
     \\ \\

    \theutterance \stepcounter{utterance}  
    & \multicolumn{4}{p{0.6\linewidth}}{
        \cellcolor[rgb]{0.561,0.737,0.561}{
            \makecell[{{p{\linewidth}}}]{
                \texttt{\tiny{[P1$\rangle$GM]}}
                \texttt{```python} \\
\texttt{SOLVE:} \\
\texttt{def f(text: str) {-}> str:} \\
\texttt{    """Returns the input string repeated three times, with '\_alt' appended if the input is a single word containing the letter 'l'."""} \\
\texttt{    if ' ' not in text and 'l' in text:} \\
\texttt{        return text * 3 + '\_alt'} \\
\texttt{    return text * 3} \\
\texttt{```} \\
            }
        }
    }
    & & \\ \\

    \theutterance \stepcounter{utterance}  
    & & & \multicolumn{2}{p{0.3\linewidth}}{
        \cellcolor[rgb]{0.9,0.9,0.9}{
            \makecell[{{p{\linewidth}}}]{
                \texttt{\tiny{[GM$|$GM]}}
                \texttt{Game Verdict: LOSS (turn limit reached)} \\
            }
        }
    }
    & & \\ \\

\end{supertabular}
}

\end{document}
